% LONG PAPER. No page limit, for a later archival venue.
%
% This is NOT a superset of neurips_2026.tex produced by copy and paste. The two papers
% are written alongside each other and share exactly one thing, numbers.tex, so that a
% corrected figure propagates to both. Prose is written separately because a five page
% argument and a full account are different documents, not one document at two lengths.
%
% What lives here and not in the workshop paper:
%   the 16 published baselines          table, section 4.1
%   the retriever comparison            table, section 3.1
%   the BM25 ablation                   table, section 3.1
%   the k sweep                         table, section 3.1
%   claim decomposition, six merges     table, section 3.1
%   the oracles                         table, section 5
%   the full cost breakdown             table, section 5
%
% Compiled with the plain article class, not the NeurIPS style, because this is not a
% NeurIPS submission and the workshop template's page furniture is wrong for it.

\documentclass[11pt]{article}

\usepackage[margin=1in]{geometry}
\usepackage[utf8]{inputenc}
\usepackage[T1]{fontenc}
\usepackage{hyperref}
\usepackage{url}
\usepackage{booktabs}
\usepackage{amsfonts}
\usepackage{nicefrac}
\usepackage{microtype}
\usepackage{xcolor}
\usepackage{natbib}

% Same macros as the workshop paper. This is the only shared file.
% numbers.tex
%
% Every number either paper cites, defined once.
% Both neurips_2026.tex (workshop, 5 pages) and long.tex \input this file,
% so the two papers cannot drift apart.
%
% Rule: a macro is added here only after the figure exists in docs/paper_numbers.md
% with an n and a reproducing script. See that file's rule 1.
% The section number in each comment points into docs/paper_numbers.md.
%
% Naming: LaTeX control sequences cannot contain digits, so 3B is ThreeB
% and 7B is SevenB throughout.


% ---------------------------------------------------------------------------
% Splits and sample sizes
% ---------------------------------------------------------------------------

\newcommand{\nDev}{700}                  % testmini, development split
\newcommand{\nTest}{1{,}700}             % test.json, held-out split
\newcommand{\nIE}{600}
\newcommand{\nMATH}{600}
\newcommand{\nKNOW}{500}


% ---------------------------------------------------------------------------
% Headline accuracy, strict scoring, n = 1,700 held out. Section 2.18
% ---------------------------------------------------------------------------

\newcommand{\accThreeB}{59.8\%}
\newcommand{\accThreeBcount}{1016/1700}
\newcommand{\accThreeBci}{57.4--62.1}

\newcommand{\accSevenB}{67.8\%}
\newcommand{\accSevenBcount}{1153/1700}
\newcommand{\accSevenBci}{65.6--70.0}

\newcommand{\accRouted}{75.8\%}
\newcommand{\accRoutedcount}{1289/1700}
\newcommand{\accRoutedci}{73.7--77.8}

\newcommand{\accCloud}{77.4\%}
\newcommand{\accCloudcount}{1315/1700}
\newcommand{\accCloudci}{75.3--79.3}

\newcommand{\gainOverSevenB}{8.0}        % routed minus 7B alone
\newcommand{\gainOverThreeB}{16.0}       % routed minus 3B alone


% ---------------------------------------------------------------------------
% Routed against cloud alone, paired McNemar on identical claims. Section 2.18
% ---------------------------------------------------------------------------

\newcommand{\mcnemarAllP}{0.119}
\newcommand{\mcnemarAllDisagree}{258}
\newcommand{\mcnemarAllRoutedWins}{116}
\newcommand{\mcnemarAllCloudWins}{142}

\newcommand{\mcnemarIEP}{0.007}          % the one significant result, routed loses
\newcommand{\mcnemarIEDisagree}{109}
\newcommand{\mcnemarMATHP}{1.000}
\newcommand{\mcnemarKNOWP}{0.917}


% ---------------------------------------------------------------------------
% Per subset, all four arms. Section 2.18
% ---------------------------------------------------------------------------

\newcommand{\accThreeBIE}{63.3\%}
\newcommand{\accThreeBMATH}{57.3\%}
\newcommand{\accThreeBKNOW}{58.4\%}

\newcommand{\accSevenBIE}{68.5\%}
\newcommand{\accSevenBMATH}{66.5\%}
\newcommand{\accSevenBKNOW}{68.6\%}

\newcommand{\accRoutedIE}{77.5\%}
\newcommand{\accRoutedMATH}{77.3\%}
\newcommand{\accRoutedKNOW}{72.0\%}

\newcommand{\accCloudIE}{82.3\%}
\newcommand{\accCloudMATH}{77.2\%}
\newcommand{\accCloudKNOW}{71.6\%}

\newcommand{\escRateIE}{36.5\%}
\newcommand{\escRateMATH}{79.8\%}
\newcommand{\escRateKNOW}{42.0\%}


% ---------------------------------------------------------------------------
% Contribution 2: same-family agreement is correlated error. Section 2.18
% The whole FDV-IE deficit sits on claims the gate kept local.
% ---------------------------------------------------------------------------

\newcommand{\keptLocalIEn}{381}
\newcommand{\keptLocalIERouted}{76.9\%}
\newcommand{\keptLocalIECloud}{85.6\%}
\newcommand{\keptLocalIEGap}{8.7}        % points, routed below cloud

\newcommand{\escalatedIEn}{219}
\newcommand{\escalatedIERouted}{78.5\%}
\newcommand{\escalatedIECloud}{76.7\%}   % routed is slightly ahead here


% ---------------------------------------------------------------------------
% Escalation earns its cost. Section 2.18
% ---------------------------------------------------------------------------

\newcommand{\cloudCalls}{908}
\newcommand{\escWouldThreeB}{45.7\%}
\newcommand{\escWouldSevenB}{60.8\%}
\newcommand{\escCloudGave}{75.8\%}
\newcommand{\escGain}{15.0}              % points over the 7B on the same claims

\newcommand{\trigDisagreeN}{459}
\newcommand{\trigDisagreeSevenB}{57.3\%}
\newcommand{\trigDisagreeCloud}{73.4\%}
\newcommand{\trigDisagreeGain}{16.1}

\newcommand{\trigNumericN}{449}
\newcommand{\trigNumericSevenB}{64.4\%}
\newcommand{\trigNumericCloud}{78.2\%}
\newcommand{\trigNumericGain}{13.8}


% ---------------------------------------------------------------------------
% Routing behaviour. Section 2.18
% ---------------------------------------------------------------------------

\newcommand{\escalationRate}{53.4\%}     % claims that called the cloud
\newcommand{\onDeviceRate}{46.6\%}       % claims answered entirely on device
\newcommand{\localsAgreeRate}{66.2\%}
\newcommand{\accKeptLocal}{75.9\%}
\newcommand{\accCloudCalled}{75.8\%}


% ---------------------------------------------------------------------------
% Cost, from our own recorded token counts. Section 2.18, and 2.18.1 and 3.5
% USD and tokens only. Never CNY.
% ---------------------------------------------------------------------------

\newcommand{\costRouted}{\$1.066}
\newcommand{\costCloud}{\$1.882}
\newcommand{\costRatio}{56.6\%}
\newcommand{\tokensRoutedIn}{3{,}400{,}854}
\newcommand{\tokensRoutedOut}{2{,}106{,}383}
\newcommand{\tokensCloudIn}{6{,}319{,}626}
\newcommand{\tokensCloudOut}{3{,}560{,}933}
\newcommand{\medianLatency}{26.0}        % seconds per claim, GPU box, report median not mean


% ---------------------------------------------------------------------------
% Oracles. Upper bounds only. Neither may be presented as a system result.
% Section 2.18
% ---------------------------------------------------------------------------

\newcommand{\oracleThreeWay}{84.8\%}
\newcommand{\oracleThreeWayci}{83.0--86.4}
\newcommand{\oracleLocalOnly}{79.9\%}    % beats the cloud with zero cloud calls
\newcommand{\oracleLocalOnlyci}{78.0--81.8}


% ---------------------------------------------------------------------------
% Unparseable output. Section 2.18
% Cloud denominator is calls made, not 1,700.
% ---------------------------------------------------------------------------

\newcommand{\unparseThreeB}{1.65\%}
\newcommand{\unparseSevenB}{2.82\%}
\newcommand{\unparseCloud}{1.65\%}
\newcommand{\unparseRouted}{0.94\%}


% ---------------------------------------------------------------------------
% Retrieval, macro recall at k = 10. Section 1 and 2.16
% These are retrieval results. They may NOT be written as accuracy drivers.
% ---------------------------------------------------------------------------

\newcommand{\recallOursDev}{74.60\%}
\newcommand{\recallOursTest}{75.16\%}
\newcommand{\recallEmbedDev}{68.01\%}    % text-embedding-3-large, best published
\newcommand{\recallEmbedTest}{69.75\%}
\newcommand{\recallUpstreamBMDev}{65.16\%}
\newcommand{\recallUpstreamBMTest}{64.29\%}
\newcommand{\recallContrieverDev}{33.48\%}
\newcommand{\recallContrieverTest}{35.02\%}

\newcommand{\recallGainDev}{6.6}         % points over best published, testmini
\newcommand{\recallGainTest}{5.4}        % points over best published, held out

\newcommand{\recallElementTest}{70.0\%}
\newcommand{\recallAllGoldTest}{49.4\%}

\newcommand{\recallOursIETest}{77.92\%}
\newcommand{\recallOursMATHTest}{80.92\%}
\newcommand{\recallOursKNOWTest}{64.94\%}

% Why ours beats upstream's BM25: an implementation gap, priced.
\newcommand{\ablationIDF}{2.51}          % Lucene IDF, points
\newcommand{\ablationPunct}{2.02}        % dropping punctuation tokens, points
\newcommand{\ablationStem}{0.06}         % stemming, worth nothing
\newcommand{\tableInflation}{1.81}       % punctuation inflates table length this much
\newcommand{\paraInflation}{1.13}        % against this for paragraphs

% Fusion is dead on both splits.
\newcommand{\recallFusionTest}{74.95\%}  % below our BM25 alone


% ---------------------------------------------------------------------------
% Contribution 4: the error taxonomy. Section 2.17
% ---------------------------------------------------------------------------

\newcommand{\missingInfoKnowDev}{63.5\%}
\newcommand{\missingInfoKnowTest}{63.4\%}
\newcommand{\missingInfoAllTest}{41.0\%}
\newcommand{\missingInfoCiteRateTest}{38.4\%}
\newcommand{\missingInfoWrongTest}{47.6\%}
\newcommand{\missingInfoBaseTest}{35.8\%}


% ---------------------------------------------------------------------------
% The refuted bias, replicated four times. Section 2.18
% Share of claims each arm answers "entailed", where gold is 50.1% entailed.
% ---------------------------------------------------------------------------

\newcommand{\saysEntailedThreeB}{39.5\%}
\newcommand{\saysEntailedSevenB}{49.2\%}  % the only nearly calibrated arm
\newcommand{\saysEntailedRouted}{41.1\%}
\newcommand{\saysEntailedCloud}{33.9\%}
\newcommand{\saysEntailedGold}{50.1\%}

\newcommand{\entailedRouted}{67.1\%}
\newcommand{\entailedCloud}{61.7\%}
\newcommand{\refutedRouted}{84.6\%}
\newcommand{\refutedCloud}{93.1\%}


% ---------------------------------------------------------------------------
% The target device. Hard constraint, not a result. CLAUDE.md and section 3.1
% ---------------------------------------------------------------------------

\newcommand{\deviceLatency}{7}           % minutes per example, 3B on the MacBook
\newcommand{\deviceSecPerToken}{0.0641}
\newcommand{\deviceRsq}{0.995}
\newcommand{\deviceMeanPrompt}{6{,}610}  % tokens


\title{Routing Financial Claim Verification Between On-Device and Cloud Models:
       A Full Account}

% Not double-blind. This version is not under anonymous review.
% TODO: fill in once the author list is settled with the industry co-authors.
\author{TODO}
\date{Draft, \today}

\begin{document}

\maketitle

\begin{abstract}
This is the extended version of a five page workshop paper. It reports the same system and
the same held-out results, and adds the material the page limit excluded: the comparison
against sixteen published baselines, the retrieval study in full, the oracle upper bounds,
the complete cost accounting, and the components we built, measured and discarded. On
\nTest{} held-out claims from FINDVER a routed pipeline running two small models on device
and escalating selectively to a frontier cloud model scores \accRouted{} against a
cloud-only baseline's \accCloud{}, a difference that is not statistically significant
(paired McNemar, \mcnemarAllDisagree{} disagreements, $p = \mcnemarAllP{}$), while using
\costRatio{} of the cloud tokens and answering \onDeviceRate{} of claims entirely on
device. We show that the aggregate parity conceals a significant loss on one subset, and
we identify its cause in the confidence signal the gate relies upon.
\end{abstract}

\tableofcontents

\section{Introduction}
\label{sec:long-intro}

% TODO. Longer than the workshop introduction and differently structured. The workshop
% version compresses the four contributions into a list; here they can each get a
% paragraph, and the deployment setting can be motivated properly rather than asserted.

\section{Related work}
\label{sec:long-related}

% TODO. Expand the workshop section. Room here for TART, ProTrix, OpenTab, GraphOTTER,
% TableRAG, FISCAL and FinVerBench individually rather than compressed into one sentence
% each, and for the leaderboard finding: FINDVER promised an evaluation platform that was
% never launched, and email submission was retired in July 2026.
%
% The supportable phrasing is "we found no other method evaluated on FINDVER" and never
% "nobody has." The citation sweep was abstract-level screening of 23 citing papers, not
% a full-text sweep. See docs/paper_numbers.md section 5.

\section{Method}
\label{sec:long-method}

% TODO. The workshop Method is compressed. Here the pipeline can be given properly:
% the loader and its data traps, the prompt versions, the extraction levels, and the
% resume logic. Section 3.1 below is the part that differs most between the two papers.

\subsection{Retrieval}
\label{sec:long-retrieval}

% The workshop paper states these results in prose across three paragraphs. Here they get
% the tables, because the retrieval study is a substantial part of the project's work and
% compressing it was the single largest cut the page limit forced.

Table~\ref{tab:retrievers} compares our BM25 implementation against the retrievers whose
rankings the benchmark ships, on both splits, using the macro recall metric the
benchmark's own evaluation script computes.

\begin{table}[ht]
\centering
\caption{Macro recall at $k{=}10$. The development column is the split every design
decision was made against; the test column was read once. Our implementation replicates
to within 0.6 points on \nTest{} claims it was never tuned on.}
\label{tab:retrievers}
\begin{tabular}{lrr}
\toprule
retriever & development, $n{=}\nDev{}$ & test, $n{=}\nTest{} $ \\
\midrule
ours, BM25                       & \recallOursDev{}        & \recallOursTest{} \\
\texttt{text-embedding-3-large}  & \recallEmbedDev{}       & \recallEmbedTest{} \\
benchmark's BM25                 & \recallUpstreamBMDev{}  & \recallUpstreamBMTest{} \\
\texttt{contriever-msmarco}      & \recallContrieverDev{}  & \recallContrieverTest{} \\
\bottomrule
\end{tabular}
\end{table}

% TODO tables still to add here, all from docs/paper_numbers.md:
%   the seven-row BM25 ablation                       section 1.2
%   the k sweep from 5 to 30                          section 1.3
%   claim decomposition, eight merge strategies       section 1.7
%   fusion including the weak-arm result              section 1.5
%   per-subset recall                                 sections 1.4 and 2.16
%
% Also belongs here and not in the workshop paper: the two methodological findings from
% section 4.5. First, any multi-query retriever must be compared against a single-query
% one at matched candidate count, since the decomposition union scored 76.14% against
% 57.88% purely by holding 30 candidates against 10. Second, a sample sized to detect a
% large effect cannot be reused to detect a small one.

\subsection{The local tier}
% TODO. Includes the model selection study the workshop paper omits entirely: qwen3:4b
% rejected on latency, qwen2.5:3b tied with the Coder variant on accuracy and lost on
% format compliance. Both from docs/paper_numbers.md sections 2.5 and 2.5.1.
% Carry the warning that every n=102 accuracy figure is inflated by an unknown amount.

\subsection{The escalation gate}
% TODO. As the workshop paper, plus the detector development history: the three Tier 1
% candidates dropped after held-out precision fell, and the Tier 2 extenders that score
% better by F1 and were still not adopted. docs/paper_numbers.md section 2.13.4.

\subsection{Why arithmetic is escalated rather than executed}
% TODO. As the workshop paper, with the full 11 August measurement table.

\section{Experimental setup}
\label{sec:long-setup}

% TODO. As the workshop paper, plus the material compressed out of it: the label
% extractor's development against 11,200 upstream response and label pairs, the extraction
% levels, and the data traps that had to be handled in the loader.

\subsection{Comparison against published baselines}
\label{sec:long-baselines}

% TODO. The sixteen-model table, cut from the workshop paper for space and restored here.
% From docs/paper_numbers.md sections 2.1 and 2.4.1.
%
% MANDATORY CAVEAT, do not drop it. Our figures are strictly scored and the published ones
% resolve unrecoverable verdicts by a coin flip. Any table mixing the two must use the
% benchmark-compatible convention, and the strict figures belong in a separate table with
% the unparseable rate beside them. The comparison is also confounded by a different
% extractor, retriever and temperature. Two different figures near 77% are in play and
% they are not the same claim: our own cloud arm at 77.0% strict on the development split
% is a controlled internal comparison, whereas the best published 2024 RAG figure of 75.0%
% is an external one.

\section{Results}
\label{sec:long-results}

% TODO. The workshop paper carries one table. This one carries all of them, from
% docs/paper_numbers.md section 2.18: the four arms with confidence intervals, the paired
% McNemar decomposition by subset, the kept-local versus escalated split that locates the
% FDV-IE deficit, the escalation value by trigger, the routing behaviour counts, the label
% bias replication, the oracles and the cost accounting.

\begin{table}[ht]
\centering
\caption{Strict accuracy on the held-out split, $n{=}\nTest{}$, with 95\% confidence
intervals. Neither oracle is implementable; both are computed with knowledge of the gold
label and bound the headroom rather than reporting a system.}
\label{tab:headline}
\begin{tabular}{lrr}
\toprule
arm & accuracy & 95\% CI \\
\midrule
3B alone                  & \accThreeB{}   & \accThreeBci{} \\
7B alone                  & \accSevenB{}   & \accSevenBci{} \\
routed                    & \accRouted{}   & \accRoutedci{} \\
cloud alone               & \accCloud{}    & \accCloudci{} \\
\midrule
\emph{oracle, two local models only} & \emph{\oracleLocalOnly{}} & \emph{\oracleLocalOnlyci{}} \\
\emph{oracle, all three arms}        & \emph{\oracleThreeWay{}}  & \emph{\oracleThreeWayci{}} \\
\bottomrule
\end{tabular}
\end{table}

\section{Limitations}
\label{sec:long-limitations}

% TODO. Everything the workshop paper compresses into half a page:
%   the tuning disclosure, stated in full
%   the FDV-IE gate weakness, found on the reported split and therefore never fixed,
%     because fixing it would forfeit the held-out claim
%   the cloud model ignoring seed and temperature, so the cloud arm is a single sample
%     rather than a reproducible measurement
%   no target-device latency for the routed pipeline
%   the detector's reliance on claim phrasing, possibly a generation artefact
%   the error taxonomy's phrase matcher, spot-checked on development data and not
%     re-read on the held-out split

\section{Conclusion}
\label{sec:long-conclusion}

% TODO.

\bibliographystyle{plainnat}
\bibliography{refs}

\end{document}
